\subsection{Integration in $k$-space}

The quasilinear diffusion coefficient (\ref{dnrf}), describing the interaction
of the electrons with a given beam $b$ at an harmonic $n$, is rewritten as
\begin{equation}
D_{b,n}^{\text{RF}}(p,\xi)=\lim_{V\rightarrow\infty}\frac{e^{2}\pi}{V}%
\iiint\frac{d^{3}k}{(2\pi)^{3}}\left|  \mathbf{E}_{b\mathbf{k}}\right|
^{2}\left|  \Theta_{\mathbf{k}}^{b,\left(  n\right)  }\right|  ^{2}%
\delta\left(  \omega_{b}-k_{\Vert}v_{\Vert}-n\Omega\right)  \label{dn}%
\end{equation}
where,
\begin{equation}
\Theta_{\mathbf{k}}^{b,\left(  n\right)  }=\frac{1}{\sqrt{2}}e_{b\mathbf{k}%
,+}e^{-i\alpha}J_{n-1}\left(  \frac{k_{\bot}v_{\bot}}{\Omega}\right)
+\frac{1}{\sqrt{2}}e_{b\mathbf{k},-}e^{+i\alpha}J_{n+1}\left(  \frac{k_{\bot
}v_{\bot}}{\Omega}\right)  +\frac{p_{\Vert}}{p_{\bot}}e_{b\mathbf{k},\Vert
}J_{n}\left(  \frac{k_{\bot}v_{\bot}}{\Omega}\right)  \label{thetak}%
\end{equation}

Here, the polarization vector in Fourier space
\begin{equation}
\mathbf{e}_{b\mathbf{k}}=\frac{\mathbf{E}_{b\mathbf{k}}}{\left|
\mathbf{E}_{b\mathbf{k}}\right|  } \label{polvec}%
\end{equation}
is introduced, whose components in the rotating field frame are
\begin{align}
e_{b\mathbf{k},+}  &  =\frac{E_{b\mathbf{k},+}}{\left|  \mathbf{E}%
_{b\mathbf{k}}\right|  }=\frac{E_{b\mathbf{k},x}+iE_{b\mathbf{k},y}}{\sqrt
{2}\left|  \mathbf{E}_{b\mathbf{k}}\right|  }\label{pol}\\
e_{b\mathbf{k},-}  &  =\frac{E_{b\mathbf{k},-}}{\left|  \mathbf{E}%
_{b\mathbf{k}}\right|  }=\frac{E_{b\mathbf{k},x}-iE_{b\mathbf{k},y}}{\sqrt
{2}\left|  \mathbf{E}_{b\mathbf{k}}\right|  }\nonumber\\
e_{b\mathbf{k},\Vert}  &  =\frac{E_{b\mathbf{k},\Vert}}{\left|  \mathbf{E}%
_{b\mathbf{k}}\right|  }=\frac{E_{b\mathbf{k},z}}{\left|  \mathbf{E}%
_{b\mathbf{k}}\right|  }\nonumber
\end{align}

The electric field associated with a given ray is described by a plane wave,
with given frequency $\omega_{b}$ and wave number $\mathbf{k}_{b}$:%

\begin{align}
\mathbf{E}_{b}\left(  \mathbf{x},t\right)   &  =\text{Re}\left[
\mathbf{E}_{b0}e^{i\left(  \mathbf{k}_{b}\cdot\mathbf{x}-\omega_{b}t\right)
}\right]  =\text{Re}\left[  \widetilde{\mathbf{E}}_{b}\right] \label{Ecomp}\\
&  =\frac{1}{2}\left[  \mathbf{E}_{b0}e^{i\left(  \mathbf{k}_{b}%
\cdot\mathbf{x}-\omega_{b}t\right)  }+\mathbf{E}_{b0}^{\ast}e^{-i\left(
\mathbf{k}_{b}\cdot\mathbf{x}-\omega_{b}t\right)  }\right]
\end{align}
so that
\begin{align}
\mathbf{E}_{b\mathbf{k}}\left(  \mathbf{k},t\right)   &  \equiv\iiint
d^{3}x\;e^{-i\mathbf{k}\cdot\mathbf{x}}\mathbf{E}_{b}\left(  \mathbf{x}%
,t\right) \nonumber\\
&  =\frac{1}{2}\left[  \mathbf{E}_{b0}e^{-i\omega_{b}t}\iiint d^{3}%
x\;e^{i\left(  \mathbf{k}_{b}-\mathbf{k}\right)  \cdot\mathbf{x}}%
+\mathbf{E}_{b0}^{\ast}e^{i\omega_{b}t}\iiint d^{3}x\;e^{i\left(
\mathbf{k}_{b}+\mathbf{k}\right)  \cdot\mathbf{x}}\right] \nonumber\\
&  =\frac{1}{2}\left[  \mathbf{E}_{b0}e^{-i\omega_{b}t}\delta\left(
\mathbf{k}_{b}-\mathbf{k}\right)  +\mathbf{E}_{b0}^{\ast}e^{i\omega_{b}%
t}\delta\left(  \mathbf{k}_{b}+\mathbf{k}\right)  \right]
\end{align}

The electric field is the sum of two contributions with different frequencies%
\begin{align}
\mathbf{E}_{b\mathbf{k}}^{+}\left(  \mathbf{k}\right)   &  =\frac{1}%
{2}\mathbf{E}_{b0}\delta\left(  \mathbf{k}_{b}-\mathbf{k}\right)  \text{ with
}\omega=\omega_{b}\\
\mathbf{E}_{b\mathbf{k}}^{-}\left(  \mathbf{k}\right)   &  =\frac{1}%
{2}\mathbf{E}_{b0}^{\ast}\delta\left(  \mathbf{k}_{b}+\mathbf{k}\right)
\text{ with }\omega=-\omega_{b}%
\end{align}

Therefore the QL diffusion tensor (\ref{dtensor}) is the sum of these two
contributions:%
\begin{equation}
\mathbb{{D}}^{\text{RF}}=\mathbb{{D}}^{\text{RF}+}+\mathbb{{D}}^{\text{RF}-}%
\end{equation}
associated with the diffusion coefficients%
\begin{align}
D_{b,n}^{\text{RF}+}(p,\xi)  &  =\lim_{V\rightarrow\infty}\frac{e^{2}\pi}%
{V}\iiint\frac{d^{3}k}{(2\pi)^{3}}\left|  \frac{1}{2}\mathbf{E}_{b0}%
\delta\left(  \mathbf{k}_{b}-\mathbf{k}\right)  \right|  ^{2}\left|
\Theta_{\mathbf{k}}^{b^{+},n}\right|  ^{2}\delta\left(  \omega_{b}-k_{\Vert
}v_{\Vert}-n\Omega\right) \label{dbnp}\\
D_{b,n}^{\text{RF}-}(p,\xi)  &  =\lim_{V\rightarrow\infty}\frac{e^{2}\pi}%
{V}\iiint\frac{d^{3}k}{(2\pi)^{3}}\left|  \frac{1}{2}\mathbf{E}_{b0}^{\ast
}\delta\left(  \mathbf{k}_{b}+\mathbf{k}\right)  \right|  ^{2}\left|
\Theta_{\mathbf{k}}^{b^{-},n}\right|  ^{2}\delta\left(  -\omega_{b}-k_{\Vert
}v_{\Vert}-n\Omega\right) \nonumber
\end{align}
with, using (\ref{thetak})%
\begin{align}
\Theta_{\mathbf{k}}^{b+,\left(  n\right)  }  &  =\frac{1}{\sqrt{2}}%
e_{b^{+}\mathbf{k},+}e^{-i\alpha}J_{n-1}\left(  \frac{k_{\bot}v_{\bot}}%
{\Omega}\right)  +\frac{1}{\sqrt{2}}e_{b^{+}\mathbf{k},-}e^{+i\alpha}%
J_{n+1}\left(  \frac{k_{\bot}v_{\bot}}{\Omega}\right)  +\frac{p_{\Vert}%
}{p_{\bot}}e_{b^{+}\mathbf{k},\Vert}J_{n}\left(  \frac{k_{\bot}v_{\bot}%
}{\Omega}\right) \label{thetakp}\\
\Theta_{\mathbf{k}}^{b-,\left(  n\right)  }  &  =\frac{1}{\sqrt{2}}%
e_{b^{-}\mathbf{k},+}e^{-i\alpha}J_{n-1}\left(  \frac{k_{\bot}v_{\bot}}%
{\Omega}\right)  +\frac{1}{\sqrt{2}}e_{b^{-}\mathbf{k},-}e^{+i\alpha}%
J_{n+1}\left(  \frac{k_{\bot}v_{\bot}}{\Omega}\right)  +\frac{p_{\Vert}%
}{p_{\bot}}e_{b^{-}\mathbf{k},\Vert}J_{n}\left(  \frac{k_{\bot}v_{\bot}%
}{\Omega}\right) \nonumber
\end{align}

The polarization components (\ref{pol}) are%
\begin{align}
e_{b^{+}\mathbf{k},+}  &  =\frac{E_{b0,x}+iE_{b0,y}}{\sqrt{2}\left|
\mathbf{E}_{b0}\right|  }=e_{b0,+}\\
e_{b^{+}\mathbf{k},-}  &  =\frac{E_{b0,x}-iE_{b0,y}}{\sqrt{2}\left|
\mathbf{E}_{b0}\right|  }=e_{b0,-}\nonumber\\
e_{b^{+}\mathbf{k},\Vert}  &  =\frac{E_{b0,z}}{\left|  \mathbf{E}_{b0}\right|
}=e_{b0,\Vert}\nonumber
\end{align}
and
\begin{align}
e_{b^{-}\mathbf{k},+}  &  =\frac{E_{b0,x}^{\ast}+iE_{b0,y}^{\ast}}{\sqrt
{2}\left|  \mathbf{E}_{b0}\right|  }=e_{b0,-}^{\ast}\\
e_{b^{-}\mathbf{k},-}  &  =\frac{E_{b0,x}^{\ast}-iE_{b0,y}^{\ast}}{\sqrt
{2}\left|  \mathbf{E}_{b0}\right|  }=e_{b0,+}^{\ast}\nonumber\\
e_{b^{-}\mathbf{k},\Vert}  &  =\frac{E_{b0,z}^{\ast}}{\left|  \mathbf{E}%
_{b0}\right|  }=e_{b0,\Vert}^{\ast}\nonumber
\end{align}

Expressing the condition $\mathbf{k}=\pm\mathbf{k}_{b}$ from the first delta
function in (\ref{dbnp}), we find%
\begin{align}
D_{b,n}^{\text{RF}+}(p,\xi)  &  =\lim_{V\rightarrow\infty}\frac{e^{2}\pi}%
{4V}\frac{1}{\left|  v_{\Vert}\right|  }\iiint\frac{d^{3}k}{(2\pi)^{3}}\left|
\mathbf{E}_{b0}\right|  ^{2}\delta^{2}\left(  \mathbf{k}_{b}-\mathbf{k}%
\right)  \left|  \Theta_{\mathbf{k}}^{b^{+},n}\right|  ^{2}\delta\left(
k_{b\Vert}-\frac{\omega_{b}-n\Omega}{v_{\Vert}}\right) \label{dbnp2}\\
D_{b,n}^{\text{RF}-}(p,\xi)  &  =\lim_{V\rightarrow\infty}\frac{e^{2}\pi}%
{4V}\frac{1}{\left|  v_{\Vert}\right|  }\iiint\frac{d^{3}k}{(2\pi)^{3}}\left|
\mathbf{E}_{b0}^{\ast}\right|  ^{2}\delta^{2}\left(  \mathbf{k}_{b}%
+\mathbf{k}\right)  \left|  \Theta_{\mathbf{k}}^{b^{-},n}\right|  ^{2}%
\delta\left(  k_{b\Vert}-\frac{\omega_{b}+n\Omega}{v_{\Vert}}\right) \nonumber
\end{align}
with now (\ref{thetakp}) being%
\begin{align}
\Theta_{\mathbf{k}}^{b+,\left(  n\right)  }  &  =\frac{1}{\sqrt{2}}%
e_{b0,+}e^{-i\alpha_{b}}J_{n-1}\left(  \frac{k_{b\bot}v_{\bot}}{\Omega
}\right)  +\frac{1}{\sqrt{2}}e_{b0,-}e^{+i\alpha_{b}}J_{n+1}\left(
\frac{k_{b\bot}v_{\bot}}{\Omega}\right)  +\frac{p_{\Vert}}{p_{\bot}%
}e_{b0,\Vert}J_{n}\left(  \frac{k_{b\bot}v_{\bot}}{\Omega}\right)
\label{thetaknp}\\
\Theta_{\mathbf{k}}^{b-,\left(  n\right)  }  &  =\frac{1}{\sqrt{2}}%
e_{b0,-}^{\ast}e^{-i\left(  \alpha_{b}+\pi\right)  }J_{n-1}\left(
\frac{k_{b\bot}v_{\bot}}{\Omega}\right)  +\frac{1}{\sqrt{2}}e_{b0,+}^{\ast
}e^{+i\left(  \alpha_{b}+\pi\right)  }J_{n+1}\left(  \frac{k_{b\bot}v_{\bot}%
}{\Omega}\right)  +\frac{p_{\Vert}}{p_{\bot}}e_{b0,\Vert}^{\ast}J_{n}\left(
\frac{k_{b\bot}v_{\bot}}{\Omega}\right) \nonumber\\
&  =-\frac{1}{\sqrt{2}}e_{b0,-}^{\ast}e^{-i\alpha_{b}}J_{n-1}\left(
\frac{k_{b\bot}v_{\bot}}{\Omega}\right)  -\frac{1}{\sqrt{2}}e_{b0,+}^{\ast
}e^{+i\alpha_{b}}J_{n+1}\left(  \frac{k_{b\bot}v_{\bot}}{\Omega}\right)
+\frac{p_{\Vert}}{p_{\bot}}e_{b0,\Vert}^{\ast}J_{n}\left(  \frac{k_{b\bot
}v_{\bot}}{\Omega}\right)  \label{thetaknm}%
\end{align}
where we used the fact that $\mathbf{k}_{b}$ is transformed into
$-\mathbf{k}_{b}$ by changing $k_{b\Vert}$ to $-k_{b\Vert}$ and $\alpha_{b}$
to $\alpha_{b}+\pi$.

Using the properties of the Bessel functions, and defining $n^{\prime}=-n$,%
\begin{align}
J_{n+1}  &  =\left(  -1\right)  ^{n+1}J_{n^{\prime}-1}\\
J_{n-1}  &  =\left(  -1\right)  ^{n-1}J_{n^{\prime}+1}\nonumber\\
J_{n}  &  =\left(  -1\right)  ^{n}J_{n^{\prime}}\nonumber
\end{align}
so that (\ref{thetaknm}) becomes%
\begin{align}
\Theta_{\mathbf{k}}^{b-,\left(  n\right)  }  &  =\left(  -1\right)
^{n}\left[  \frac{1}{\sqrt{2}}e_{b0,-}^{\ast}e^{-i\alpha_{b}}J_{n^{\prime}%
+1}\left(  \frac{k_{b\bot}v_{\bot}}{\Omega}\right)  +\frac{1}{\sqrt{2}%
}e_{b0,+}^{\ast}e^{+i\alpha_{b}}J_{n^{\prime}-1}\left(  \frac{k_{b\bot}%
v_{\bot}}{\Omega}\right)  +\frac{p_{\Vert}}{p_{\bot}}e_{b0,\Vert}^{\ast
}J_{n^{\prime}}\left(  \frac{k_{b\bot}v_{\bot}}{\Omega}\right)  \right] \\
&  =\left(  -1\right)  ^{n}\left(  \Theta_{\mathbf{k}}^{b+,\left(  n\right)
^{\prime}}\right)  ^{\ast}%
\end{align}
and the diffusion coefficients become%
\begin{align}
D_{b,n}^{\text{RF}+}(p,\xi)  &  =\lim_{V\rightarrow\infty}\frac{e^{2}\pi}%
{4V}\frac{1}{\left|  v_{\Vert}\right|  }\left|  \Theta_{\mathbf{k}}^{b,\left(
n\right)  }\right|  ^{2}\left|  \mathbf{E}_{b0}\right|  ^{2}\delta\left(
k_{b\Vert}-\frac{\omega_{b}-n\Omega}{v_{\Vert}}\right)  \iiint\frac{d^{3}%
k}{(2\pi)^{3}}\delta^{2}\left(  \mathbf{k}_{b}-\mathbf{k}\right) \\
D_{b,n}^{\text{RF}-}(p,\xi)  &  =\lim_{V\rightarrow\infty}\frac{e^{2}\pi}%
{4V}\frac{1}{\left|  v_{\Vert}\right|  }\left|  \Theta_{\mathbf{k}}^{b,\left(
n\right)  ^{\prime}}\right|  ^{2}\left|  \mathbf{E}_{b0}\right|  ^{2}%
\delta\left(  k_{b\Vert}-\frac{\omega_{b}-n^{\prime}\Omega}{v_{\Vert}}\right)
\iiint\frac{d^{3}k}{(2\pi)^{3}}\delta^{2}\left(  \mathbf{k}_{b}+\mathbf{k}%
\right) \nonumber
\end{align}

Using Parseval's theorem,
\begin{align}
\lim_{V\rightarrow\infty}\frac{1}{V}\iiint\frac{d^{3}k}{(2\pi)^{3}}\delta
^{2}\left(  \mathbf{k}_{b}-\mathbf{k}\right)   &  =\lim_{V\rightarrow\infty
}\frac{1}{V}\iiint d^{3}x=1\nonumber\\
\lim_{V\rightarrow\infty}\frac{1}{V}\iiint\frac{d^{3}k}{(2\pi)^{3}}\delta
^{2}\left(  \mathbf{k}_{b}+\mathbf{k}\right)   &  =\lim_{V\rightarrow\infty
}\frac{1}{V}\iiint d^{3}x=1
\end{align}
one obtains
\begin{align}
D_{b,n}^{\text{RF}+}(p,\xi)  &  =e^{2}\pi\frac{1}{\left|  v_{\Vert}\right|
}\frac{\left|  \mathbf{E}_{b0}\right|  ^{2}}{4}\left|  \Theta_{\mathbf{k}%
}^{b,\left(  n\right)  }\right|  ^{2}\delta\left(  k_{b\Vert}-\frac{\omega
_{b}-n\Omega}{v_{\Vert}}\right) \\
D_{b,n}^{\text{RF}-}(p,\xi)  &  =e^{2}\pi\frac{1}{\left|  v_{\Vert}\right|
}\frac{\left|  \mathbf{E}_{b0}\right|  ^{2}}{4}\left|  \Theta_{\mathbf{k}%
}^{b,\left(  n\right)  ^{\prime}}\right|  ^{2}\delta\left(  k_{b\Vert}%
-\frac{\omega_{b}-n^{\prime}\Omega}{v_{\Vert}}\right)
\end{align}

In the expression (\ref{dtensor2}-\ref{dtensor4}) for $\mathbb{{D}}%
^{\text{RF}-}$ tensor elements, we subsitute%
\begin{equation}
\left[  1-\xi^{2}-\frac{n\Omega}{-\omega_{b}}\right]  =\left[  1-\xi^{2}%
-\frac{n^{\prime}\Omega}{\omega_{b}}\right]
\end{equation}

Then, by redefining $n^{\prime}\rightarrow n$ in the sum over all harmonics
for $\mathbb{{D}}^{\text{RF}-}$, we can combine the two contributions of
$\mathbb{{D}}^{\text{RF}+}$ and $\mathbb{{D}}^{\text{RF}-}$, and finally we
obtain an expression with one diffusion coefficient:\qquad\qquad\qquad\qquad%
\begin{align}
D_{b,n}^{\text{RF}}(p,\xi) &  =e^{2}\pi\frac{1}{\left|  v_{\Vert}\right|
}\frac{\left|  \mathbf{E}_{b0}\right|  ^{2}}{2}\left|  \Theta_{\mathbf{k}%
}^{b,\left(  n\right)  }\right|  ^{2}\delta\left(  k_{b\Vert}-\frac{\omega
_{b}-n\Omega}{v_{\Vert}}\right)  \nonumber\\
&  =e^{2}\pi\frac{1}{\left|  v_{\Vert}\right|  }\frac{\left|  \mathbf{E}%
_{b0}\right|  ^{2}}{2}\left|  \Theta_{\mathbf{k}}^{b,\left(  n\right)
}\right|  ^{2}\frac{c}{\omega_{b}}\delta\left(  N_{b\Vert}-N_{\Vert\text{res}%
}\right)  \label{diffcoef}%
\end{align}
with (\ref{thetaknp})%
\begin{equation}
\Theta_{\mathbf{k}}^{b,\left(  n\right)  }=\frac{1}{\sqrt{2}}e_{b0,+}%
e^{-i\alpha_{b}}J_{n-1}\left(  \frac{k_{b\bot}v_{\bot}}{\Omega}\right)
+\frac{1}{\sqrt{2}}e_{b0,-}e^{+i\alpha_{b}}J_{n+1}\left(  \frac{k_{b\bot
}v_{\bot}}{\Omega}\right)  +\frac{p_{\Vert}}{p_{\bot}}e_{b0,\Vert}J_{n}\left(
\frac{k_{b\bot}v_{\bot}}{\Omega}\right)  \label{thetakn}%
\end{equation}%
\begin{align}
N_{b\Vert} &  =\frac{k_{b\Vert}c}{\omega_{b}}\\
N_{\Vert\text{res}} &  =\frac{c}{v_{\Vert}}\left(  1-\frac{n\Omega}{\omega
_{b}}\right)  \nonumber\\
&  =\frac{1}{\beta_{Te}}\frac{p_{Te}}{p_{\Vert}}\left(  \gamma-\frac
{n^{\prime}\omega_{ce}}{\omega_{b}}\right)  \label{nparres}%
\end{align}
where
\begin{align}
\Omega &  =\frac{q_{e}B}{\gamma m_{e}}=\frac{-\omega_{ce}}{\gamma}\text{ with
}\omega_{ce}=\frac{eB}{m_{e}}\\
n^{\prime} &  =-n\\
\beta_{Te} &  =\frac{p_{Te}}{m_{e}c}=\frac{v_{Te}}{c}=\sqrt{\frac{kT_{e}%
}{m_{e}c^{2}}}%
\end{align}

\subsection{Incident Energy Flow Density}

\subsubsection{Relation to the Electric Field}

The time-averaged energy flow density in the beam is in general given by
\begin{equation}
\overline{s}_{b}=\mathbf{P}_{b}+\mathbf{T}_{b} \label{sb}%
\end{equation}
where
\begin{equation}
\mathbf{P}_{b}=\frac{1}{2}\text{Re}\left[  \widetilde{\mathbf{E}}_{b}%
\times\widetilde{\mathbf{H}}_{b}^{*}\right]
\end{equation}
is the flow of the electromagnetic energy or Poynting vector, and
\begin{equation}
\mathbf{T}_{b}=-\frac{\epsilon_{0}\omega_{b}}{4}\widetilde{\mathbf{E}}_{b}%
^{*}\cdot\frac{\partial\mathbb{{K}^{H}}}{\partial\mathbf{k}}\cdot
\widetilde{\mathbf{E}}_{b}%
\end{equation}
is the kinetic energy flow density where $\mathbb{{K}^{H}}$ is the Hermitian
part of the dielectric tensor.

$\widetilde{\mathbf{E}}_{b}$ is the complex form of the electric field
(\ref{Ecomp})
\begin{equation}
\widetilde{\mathbf{E}}_{b}=\mathbf{E}_{b0}e^{i\mathbf{k}_{b}\mathbf{\cdot
x}-i\omega_{b}t}%
\end{equation}
and the magnetic field $\widetilde{\mathbf{H}}_{b}$ is given by
\begin{equation}
\widetilde{\mathbf{H}}_{b}=\frac{\mathbf{k}_{b}\times\widetilde{\mathbf{E}%
}_{b}}{\mu_{0}\omega_{b}}=\frac{\mathbf{N}_{b}\times\widetilde{\mathbf{E}}%
_{b}}{\mu_{0}c} \label{Hfield}%
\end{equation}
using the relation $\mathbf{N}_{b}=\mathbf{k}_{b}\omega_{b}/c.$

The energy flow density (\ref{sb}) can be formally rewritten as
\begin{equation}
\overline{s}_{b}=\frac{\epsilon_{0}c}{2}\left|  \mathbf{E}_{b0}\right|
^{2}\mathbf{\Phi}_{b}%
\end{equation}
where $\mathbf{\Phi}_{b}$ is a non dimensional vector defined as
\begin{equation}
\mathbf{\Phi}_{b}=\mathbf{\Phi}_{bP}+\mathbf{\Phi}_{bT} \label{phib2}%
\end{equation}
with
\begin{equation}
\mathbf{\Phi}_{bP}=\frac{\mu_{0}c}{\left|  \mathbf{E}_{b0}\right|  ^{2}%
}\text{Re}\left[  \widetilde{\mathbf{E}}_{b}\times\widetilde{\mathbf{H}}%
_{b}^{\ast}\right]  \label{phib}%
\end{equation}
and
\begin{equation}
\mathbf{\Phi}_{bT}=\mathbf{-}\frac{1}{2}\mathbf{e}_{b}^{\ast}\cdot
\frac{\partial\mathbb{{K}^{H}}}{\partial\mathbf{N}}\cdot\mathbf{e}_{b}
\label{phit}%
\end{equation}

Using (\ref{Hfield}), we obtain
\begin{align}
\mathbf{\Phi}_{bP}  &  =\text{Re}\left[  \mathbf{e}_{b}\times\left(
\mathbf{N}_{b}\times\mathbf{e}_{b}^{\ast}\right)  \right] \nonumber\\
&  =\text{Re}\left[  \mathbf{N}_{b}-\left(  \mathbf{N}_{b}\cdot\mathbf{e}%
_{b}\right)  \mathbf{e}_{b}^{\ast}\right]  \label{phibp2}%
\end{align}

In vacuum, $\mathbf{\Phi}_{b}=\widehat{k}$ , unit vector in the direction of
propagation, and $\mathbf{\Phi}_{bT}=0$. Here,
\begin{equation}
\mathbf{e}_{b}=\frac{\mathbf{E}_{b0}}{\left|  \mathbf{E}_{b0}\right|  }%
\end{equation}
is the polarization vector. The incident power flux on a flux-surface is
therefore given by
\begin{equation}
\overline{\mathbf{s}}_{b\text{inc}}=\frac{\epsilon_{0}c}{2}\left|
\mathbf{E}_{b0}\right|  ^{2}\left|  \mathbf{\Phi}_{b}\cdot\widehat{\psi
}\right|
\end{equation}
where $\widehat{\psi}$ is the\ local vector normal to the flux-surface.

\subsubsection{Relation to the ray trajectories}

\bigskip The energy flows in the direction of the group velocity, noted
$\widehat{g}$%
\begin{equation}
\widehat{g}_{b}=\frac{\mathbf{\Phi}_{b}}{\left|  \mathbf{\Phi}_{b}\right|  }%
\end{equation}
In ray-tracing calculations, this is the direction of the ray trajectories,
which are parametrized by $\left(  \psi_{b}\left(  s\right)  ,\theta
_{b}\left(  s\right)  ,\phi_{b}\left(  s\right)  \right)  $ where $s$ is the
distance along the ray. Therefore $\widehat{g}_{b}$ can be determined from
ray-tracing calculations or any other wave propagation model, and we have
\begin{equation}
\overline{\mathbf{s}}_{b\text{inc}}=\frac{\epsilon_{0}c}{2}\left|
\mathbf{E}_{b0}\right|  ^{2}\left|  \mathbf{\Phi}_{b}\right|  \left|
\widehat{g}_{b}\cdot\widehat{\psi}\right|
\end{equation}

In DKE calculations, flux-surfaces are discretized on a grid $\psi_{l+1/2}$
and have a finite volume (). In consequence, the factor $\left|  \widehat
{g}_{b}\cdot\widehat{\psi}\right|  $ must be ''integrated'' along the ray
trajectory within this volume. We therefore define an incidence factor
\begin{equation}
f_{\text{inc},b}^{l+1/2}=\frac{1}{\Delta r_{l+1/2}^{\theta_{b}}}\int
_{\psi_{l+1/2}\mp\Delta\psi_{l+1/2}/2}^{\psi_{l+1/2}\pm\Delta\psi_{l+1/2}/2}ds
\end{equation}
where the $\pm$ sign accounts for the fact that the ray may be locally
directed either inward or outward, and where $\Delta r_{l+1/2}^{\theta_{b}}$
is the local thickness of the flux-surface.
\begin{equation}
\Delta r_{l+1/2}^{\theta_{b}}=r(\psi_{l+1,}\theta_{b})-r(\psi_{l,}\theta_{b})
\end{equation}

The factor $f_{\text{inc},b}^{l+1/2}$ can be numerically calculated from the
equilibrium and ray-tracing data. In the limit of an infinitely thin
flux-surface, we have
\begin{equation}
f_{\text{inc},b}^{l+1/2}\rightarrow\left|  \widehat{g}_{b}\cdot\widehat{\psi
}\right|  ^{-1}%
\end{equation}

We therefore substitute%
\begin{equation}
\overline{\mathbf{s}}_{b\text{inc}}=\frac{\epsilon_{0}c}{2}\left|
\mathbf{E}_{b0}\right|  ^{2}\frac{\left|  \mathbf{\Phi}_{b}\right|
}{f_{\text{inc},b}^{l+1/2}} \label{sbinc}%
\end{equation}

\subsection{Narrow Beam Approximation}

Let consider a ray which crosses the flux surface at the respective poloidal
locations $\theta_{b}$. Here, it is assumed that all electrons interact with
the beam, whatever its toroidal location, because of the axisymmetry.
Therefore, calculations are only valid for circulating electrons, except for
those localized on low $n$ order rational $q$-surfaces, since their
trajectories are exactly periodic in configuration space. Consequently, either
they may never crosses the beam or if they cross it, the quasilinear
assumption might likely fail, since some correlations could remain between two
momentum kicks. However, the total weight of these electrons is marginal, and
their contribution to RF power absorption and current generation is simply
neglected. It is important to recal that trapped particle interaction with RF
wave is not addressed in this code for similar arguments. For all particles
whose trajectory is periodic, a specific treatment of the quasilinear
interaction is necessary, which is at this stage beyond the scope of this development.

If we assume that the extension of the beam power on the flux-surface is very
small in the poloidal directions, we can approximate
\begin{equation}
\overline{\mathbf{s}}_{b\text{inc}}\simeq\frac{P_{b,\text{inc}}}%
{\mathcal{A}_{b}}\delta\left(  \theta-\theta_{b}\right)
\end{equation}

This assumption is in contradication with the plane-wave decomposition of the
RF beam, since plane waves are essentially present in the entire configuration
space. However, it is assumed that plane waves interfere destructively
everywhere except in the beam region, which can be limited to a narrow
poloidal and toroidal extention. By definition, $\mathcal{A}_{b}$ is
determined so that $P_{b,\text{inc}}$ is the total incident power on the
flux-surface
\begin{equation}
\iint dS\;\overline{\mathbf{s}}_{b\text{inc}}=P_{b,\text{inc}}%
\end{equation}
which gives in the system $\left(  \psi,\theta,\phi\right)  $ , using
(\ref{dsdthetadphi})
\begin{equation}
1=\int_{0}^{2\pi}d\phi\;\int_{0}^{2\pi}\frac{rR}{\mathcal{A}_{b}\left|
\widehat{\psi}\cdot\widehat{r}\right|  }\delta\left(  \theta-\theta
_{b}\right)  d\theta
\end{equation}

Hence,
\begin{equation}
\mathcal{A}_{b}=\frac{2\pi rR}{\left|  \widehat{\psi}\cdot\widehat{r}\right|
}%
\end{equation}
and the incident power flux is simply%

\begin{equation}
\overline{\mathbf{s}}_{b\text{inc}}=\frac{P_{b,\text{inc}}}{2\pi rR}\left|
\widehat{\psi}\cdot\widehat{r}\right|  \delta\left(  \theta-\theta_{b}\right)
\label{sbinc2}%
\end{equation}

Using (\ref{sbinc}) and (\ref{sbinc2}),
\begin{equation}
\frac{\left|  \mathbf{E}_{b0}\right|  ^{2}}{2}=\frac{1}{\epsilon_{0}c}\frac
{1}{2\pi rR}\left|  \widehat{\psi}\cdot\widehat{r}\right|  \frac
{f_{\text{inc},b}^{l+1/2}}{\left|  \mathbf{\Phi}_{b}\right|  }\delta\left(
\theta-\theta_{b}\right)  P_{b,\text{inc}} \label{e0}%
\end{equation}

For concentric circular flux surfaces, $\widehat{\psi}=\widehat{r}$ , then
$\left|  \widehat{\psi}\cdot\widehat{r}\right|  =1$. We obtain
\begin{equation}
\frac{\left|  \mathbf{E}_{b0}\right|  ^{2}}{2}=\frac{1}{\epsilon_{0}c}\frac
{1}{2\pi rR_{p}}\frac{\Psi}{\Psi_{p}}\frac{f_{\text{inc},b}^{l+1/2}}{\left|
\mathbf{\Phi}_{b}\right|  }\delta\left(  \theta-\theta_{b}\right)  P_{b,\text{
inc}}%
\end{equation}
where $\Psi_{p}\left(  r\right)  =\Psi_{p}\left(  r,\pi/2\right)  =1+\epsilon$.

\subsection{Normalized Diffusion Coefficient}

The diffusion coefficient is usually normalized to the collisional diffusion
coefficient $\nu_{e}p_{Te}^{2}$, thus defining
\begin{equation}
\overline{D}_{b,n}^{\text{RF}}(p,\xi)\equiv\frac{D_{b,n}^{\text{RF}}(p,\xi
)}{\nu_{e}p_{Te}^{2}} \label{difcoeffnorm}%
\end{equation}
with
\begin{equation}
\nu_{e}=\frac{e^{4}n_{e}\ln\Lambda}{4\pi\epsilon_{0}^{2}m_{e}^{2}v_{Te}^{3}}%
\end{equation}

Using (\ref{diffcoef}) and (\ref{e0}), the diffusion coefficient
(\ref{difcoeffnorm}) becomes
\begin{equation}
\overline{D}_{b,n}^{\text{RF}}(p,\xi)=\overline{D}_{b,n,0}^{\text{RF}}\left|
\Theta_{\mathbf{k}}^{b,\left(  n\right)  }\right|  ^{2}\frac{\gamma p_{Te}%
}{p\left|  \xi\right|  }\delta\left(  N_{b\Vert}-N_{\Vert\text{res}}\right)
2\pi\delta\left(  \theta-\theta_{b}\right)  \label{difcoeffnorm2}%
\end{equation}
with
\begin{align}
\overline{D}_{b,n,0}^{\text{RF}} &  =\frac{1}{\nu_{e}p_{Te}^{2}}\frac
{1}{v_{Te}}\frac{\pi}{\omega_{b}}\frac{e^{2}}{\epsilon_{0}}\frac{1}{4\pi
^{2}rR}\left|  \widehat{\psi}\cdot\widehat{r}\right|  \frac{f_{\text{inc}%
,b}^{l+1/2}}{\left|  \mathbf{\Phi}_{b}\right|  }P_{b,\text{inc}}\nonumber\\
&  =\frac{1}{rR}\frac{1}{m_{e}\ln\Lambda}\frac{1}{\omega_{b}\omega_{pe}^{2}%
}\left|  \widehat{\psi}\cdot\widehat{r}\right|  \frac{f_{\text{inc},b}%
^{l+1/2}}{\left|  \mathbf{\Phi}_{b}\right|  }P_{b,\text{inc}}\label{drf0}%
\end{align}
where $\omega_{pe}$ is the electron plasma frequency
\begin{equation}
\omega_{pe}=\sqrt{\frac{e^{2}n_{e}}{\epsilon_{0}m_{e}}}%
\end{equation}

We also recall the expressions (\ref{thetakn}), (\ref{nparres}) and
(\ref{phib2}-\ref{phibp2})%
\begin{equation}
\Theta_{\mathbf{k}}^{b,\left(  n\right)  }=\frac{1}{\sqrt{2}}e_{b0,+}%
e^{-i\alpha_{b}}J_{n-1}\left(  \frac{k_{b\bot}v_{\bot}}{\Omega}\right)
+\frac{1}{\sqrt{2}}e_{b0,-}e^{+i\alpha_{b}}J_{n+1}\left(  \frac{k_{b\bot
}v_{\bot}}{\Omega}\right)  +\frac{p_{\Vert}}{p_{\bot}}e_{b0,\Vert}J_{n}\left(
\frac{k_{b\bot}v_{\bot}}{\Omega}\right)  \label{thetakn2}%
\end{equation}%
\begin{equation}
N_{\Vert\text{res}}=\frac{c}{v_{\Vert}}\left(  1-\frac{n\Omega}{\omega_{b}%
}\right)  =\frac{1}{\beta_{Te}}\frac{p_{Te}}{p_{\Vert}}\left(  \gamma
-\frac{n^{\prime}\omega_{ce}}{\omega_{b}}\right)  \label{nparres2}%
\end{equation}%
\begin{align}
\mathbf{\Phi}_{b}  &  =\mathbf{\Phi}_{bP}+\mathbf{\Phi}_{bT}\label{phib3}\\
\mathbf{\Phi}_{bP}  &  =\text{Re}\left[  \mathbf{N}_{b}-\left(  \mathbf{N}%
_{b}\cdot\mathbf{e}_{b}\right)  \mathbf{e}_{b}^{\ast}\right] \label{phibp3}\\
\mathbf{\Phi}_{bT}  &  =\mathbf{-}\frac{1}{2}\mathbf{e}_{b}^{\ast}\cdot
\frac{\partial\mathbb{{K}^{H}}}{\partial\mathbf{N}}\cdot\mathbf{e}_{b}
\label{phibt3}%
\end{align}

For concentric circular flux surfaces,
\begin{equation}
\overline{D}_{b,n,0}^{\text{RF}}=\frac{1}{rR_{p}}\frac{\Psi}{\Psi_{p}}\frac
{1}{m_{e}\ln\Lambda}\frac{1}{\omega_{b}\omega_{pe}^{2}}\frac{f_{\text{inc}%
,b}^{l+1/2}}{\left|  \mathbf{\Phi}_{b}\right|  }P_{b,\text{inc}}%
\end{equation}
